\chapter{Összefoglaló}
\label{ch:summary}

A dolgozat célja a fotó osztályozási probléma vizsgálata volt neurális hálózatokkal saját nagyméretű címkézett adatbázison. 
A vizsgálat során összegyűjtöttem 180750 egyedi fotót az elmúlt 20 év saját készítésű képeiből, ebből címkézett adatbázist készítettem SBphotothemes180k néven kétféle címkézéssel.
Az ImageNet1K címkékből 114 kategória adta az első osztályozást a 300 képnél kisebb osztályok megszüntetése után, a Places365 címkékből 121 kategória adta a második osztályozást. Ezek után 80-10-10 arányban felosztottam tanító, validáló és tesztelő adathalmazra az adatbázist. Az adathalmaz képeit 224*244 méretű négyzetes jpg 80 tömörítésű képekké konvertáltam. Így a közel 2TB mennyiségű forrásképből nagyjából 10GB méretű adatbázis készült.
Az elkészült adatbázison vizsgáltam az osztályozási eloszlásokat, és az EXIF adatokat. Az osztályozások estén az átlagos osztályméret 1500 kép körül alakult, a medián 1000 kép körül alakult ami jó, de mivel természetes adatokról beszélünk, így az aszimmetria magas, 8.4 és a szórás is magas 3200 körüli tehát tipikus Long-tail eloszlást mutat.
Az EXIF adatok konverzió közbeni megőrzése és elemzése különösen fontos volt a részletes statisztikák elvégzéséhez és a későbbi vizsgálatok lehetőségeinek biztosításához. 
Az így elkészített adatbázis segítségével három fázisos vizsgálatot végeztem hét ismert és gyakran használt modellen. Ezek a modellek a ResNet50, DenseNet121, EfficientNet-B0, ConvNeXt-Tiny, ViT-B/16, Inception-v3, NASNet voltak. 
Az alkalmazott metrikák a következőek lettek: Top-1 és Top-5 pontosság, Precision, Recall, F1-score. Ezeken túl vizsgáltam a tanítás közbeni validációs és tanítási pontosságokat.
Mindhárom vizsgálatot ugyanazon az teszt adathalmazon végeztem, de más-más tanítóhalmazzal tanítottam. A három vizsgálati fázis a tanítás nélküli tesztelés, majd a teljes tanítás utáni vizsgálat volt, végül az iteratívan növekvő tanítóhalmaz vizsgálata következett. 
A vizsgálat első lépése a tanítás nélküli modellek tesztelése volt a beépített súlyokkal. Ekkor a modellek 50\% körüli Top 1 és 80\% körüli Top 5 pontosságot és 40\% körüli Precision-t nyújtottak. Itt a ConvNeXt-Tiny és ViT-B/16 emelkedtek ki.
A második fázis során a teljes 144000 képes tanítóhalmazon tanítottam minden modellt 10 epoch-on keresztül. Ekkor a modellek 75\% körüli Top 1 és 90\% körüli Top 5 pontosságot és 70\% körüli Precision-t nyújtottak a EfficientNet-B0 és NASNet emelkedtek ki.
Az utolsó vizsgálat során mindig a tanítatlan modellről indulva 1000, 5000, 10000, 25000 képes tanítóhalmazon vizsgáltam a modelleket szintén 10 epoch-on keresztül.
Az adataink alapján a ConvNeXt-Tiny és az EfficientNet-B0 modellek kiemelkedtek 25000 mintaszámmal történő tanítás után, míg a ResNet50 és DenseNet121 modellek gyengébben teljesítettek az alacsonyabb mintaszámokon, de növekvő mintaszámmal jelentős javulást mutattak. Eközben az Inception-v3 és a ViT-B/16 pedig a kis adathalmazon mutatott pontosság terén emelkedtek ki.
Ezután megkíséreltem elkészíteni egy kis méretű, lokális alkalmazásra optimalizált erőforrástakarékos modellt a ResNet18-ra építve egy saját fejrész kifejlesztésével. 
Ennek során 4 modellt készítettem, egy egyrétegű, egy mély, és két modern fejrésszel rendelkezőt. Ezen a négy modellen is elvégeztem a fent ismertetett három mérési fázist. Az egyrétegű fej meglepően jól teljesített, a mély fejrész viszont jelentősen alulmaradt a fenti metrikákban, a modern fejrész csak kicsit nyújtott nagyobb pontosságot, mint az egyrétegű, de jelentősen gyorsabban futott. Összességében a kifejlesztett modell csak kis mértékben maradt alul a vizsgált hatalmas modellektől, de futási teljesítménye sokkal jobb azokénál így nem pontosságkritikus, de erőforráskímélő megoldásokhoz, például hordozható eszközök számára ideális.
Összességében a nagyméretű adathalmaz, a többféle osztályozás, a nagy mennyiségű mért adat és a saját modellek további vizsgálatok alapját képezhetik a jövőben.